% !TEX program = lualatex
%% ============================================================
%% NRO-PRO-014-1: General onboarding: projects practical technologies
%% introduction (a apostila NeuroDev, a formação técnica da equipe)
%%
%% O conteúdo de cada parte fica em partes/. Para compilar:
%%   latexmk -lualatex NRO-PRO-014-1.tex     (ou, no Overleaf, o compilador LuaLaTeX)
%% ============================================================
\documentclass{nro-apostila}

\documento{
  codigo        = NRO-PRO-014,
  numero        = 1,
  revisao       = A,
  emissor       = {Departamento de Pesquisa e Desenvolvimento},
  tipo          = {Manual ou guia},
  categoria     = {Apostila de formação técnica},
  titulo        = {General onboarding: projects practical technologies introduction},
  subtitulo     = {NeuroDev: hardware, firmware e software para dispositivos médicos e sistemas embarcados, do básico ao avançado.},
  eixos         = {Hardware\hspace{9mm}Firmware\hspace{9mm}Software\hspace{9mm}Integração},
  data          = {25 de outubro de 2026},
  redator       = {Gabriel Matos Veneziano},
  rotuloproduto = {No Connect},
  autoria       = {NeuroDynamics},
}

\begin{document}

\capa
\sumario

%% ------------------------------------------------------------
%% INTRODUÇÃO
%% ------------------------------------------------------------
\pagina{Introdução}{Boas-vindas à NeuroDev}
\lede{A NeuroDev é o caminho de formação técnica da equipe NeuroDynamics. Ela leva qualquer membro, do primeiro dia até o nível de quem projeta e revisa sistemas completos, por quatro áreas: hardware, firmware, software e integração.}

Os conceitos são gerais e valem para qualquer sistema embarcado ou dispositivo médico. Para que o estudo não fique no abstrato, cada tópico mostra onde aquele conceito aparece no \textbf{Connect}, a plataforma de eletroestimulação que a equipe desenvolve. Assim, quem estuda aprende a teoria e, ao mesmo tempo, conhece o produto por dentro.

A apostila atende a dois públicos. Para quem está chegando, é um curso: comece pela Parte~0 e siga a trilha da sua área. Para quem já é da equipe, é uma referência: use o sumário para ir direto ao tópico e às fontes de estudo.

\begin{estatisticas}[4]
  \estatistica{\totalpartes}{Partes}
  \estatistica{\totaltopicos}{Tópicos}
  \estatistica{3}{Níveis por trilha}
  \estatistica{\totalexercicios}{Exercícios}
\end{estatisticas}

\secao{Princípios da NeuroDev}
\begin{cartoes}[2]
  \cartao{Geral, com os pés no produto}{Cada conceito é ensinado de forma geral e logo depois localizado no Connect.}
  \cartao{Do porquê ao como}{Antes de implementar, entenda por que a solução existe e que problema ela resolve.}
  \cartao{Aprender fazendo}{Todo tópico termina em exercícios práticos, de preferência no código e no esquemático reais.}
  \cartao{Documento vivo}{A apostila melhora a cada atualização. Quem estuda um tópico deixa-o melhor para o próximo.}
\end{cartoes}

%% ------------------------------------------------------------
\pagina{Guia}{Como usar esta apostila}
A NeuroDev tem cinco partes. A Parte~0 é obrigatória para todos. As Partes~1, 2 e 3 são trilhas de área, cada uma com três níveis: básico, intermediário e avançado. A Parte~4 fecha a formação e é comum a todas as áreas.

Também é altamente recomendável que, ao estudar um tópico, você crie o hábito de pesquisar o assunto na internet e de estudar por conta própria, orientado pelas suas necessidades. Guarde as fontes estudadas e contribua com a apostila! Nem todos temos a mesma bagagem de conhecimento: a apostila serve como guia e exposição dos tópicos mais importantes, mas o verdadeiro aprendizado vem da curiosidade, da descoberta própria e da resolução de problemas. Seja consciente e sincero consigo mesmo ao avançar nos tópicos.

\par\vspace{2mm}
\begin{tabularx}{\linewidth}{@{}L{14mm}L{50mm}L{40mm}Y@{}}\reguatopo
  \cab{Parte} & \cab{Conteúdo} & \cab{Para quem} & \cab{Tópicos} \\ \reguameio
  0 & \textbf{Fundamentos comuns}   & Todos, no início    & \cod{0.1} a \cod{0.7} (7) \\ \reguafina
  1 & \textbf{Hardware}             & Trilha de hardware  & \cod{1.1} a \cod{1.19} (19) \\ \reguafina
  2 & \textbf{Firmware}             & Trilha de firmware  & \cod{2.1} a \cod{2.29} (29) \\ \reguafina
  3 & \textbf{Software}             & Trilha de software  & \cod{3.1} a \cod{3.24} (24) \\ \reguafina
  4 & \textbf{Integração do sistema} & Todos, ao final    & \cod{4.1} a \cod{4.6} (6) \\ \reguabase
\end{tabularx}

\secao{Anatomia de um tópico}
Todos os tópicos seguem o mesmo formato, para que você saiba onde encontrar cada coisa:

\par\vspace{1mm}
\begin{tabularx}{\linewidth}{@{}>{\fonttitulo\color{nro-axon}}L{36mm}Y@{}}\reguatopo
  1.7 Título & Número e nome do tópico. O número diz a parte e a ordem sugerida. \\ \reguafina
  Texto de abertura & O conceito explicado de forma geral, válido para qualquer projeto. \\ \reguafina
  \rotulo{No Connect} & Quadro que mostra onde o conceito aparece nos repositórios e esquemáticos da equipe. \\ \reguafina
  \rotulo{Para estudar} & Sites, livros, vídeos e documentos internos. Os itens sublinhados são links. \\ \reguafina
  \rotulo{Pratique} & Exercícios para fixar, quase sempre ligados ao código ou ao hardware real. \\ \reguabase
\end{tabularx}

\secao{Como estudar bem}
\begin{enumerate}
  \item Leia o texto de abertura e explique com suas palavras por que aquilo existe.
  \item Use as referências para responder aos exercícios. Não é preciso ler tudo de ponta a ponta.
  \item Tente os exercícios antes de abrir o código ou pedir ajuda; anote onde travou.
  \item Leve o tópico para o repositório: encontre no código ou no esquemático o que o quadro No Connect descreve.
  \item Marque o tópico na Ficha de progresso, no fim da apostila.
  \item E principalmente, anote suas dúvidas e converse com seus colegas de equipe sobre elas. Discuta soluções, resolva problemas e não deixe suas incertezas quanto aos tópicos acumularem.
\end{enumerate}

%% ------------------------------------------------------------
\pagina{Roteiro}{Mapa das trilhas}
Todo membro faz a Parte~0, segue de básico a avançado na trilha da sua área e termina na Parte~4. As linhas tracejadas são pontes: tópicos de trilhas diferentes que tratam do mesmo assunto, cada um do seu lado.

\par\vspace{3mm}\noindent
\begin{tikzpicture}[x=1mm, y=-1mm]
  \node[nro bloco, minimum width=166mm, minimum height=11mm, align=center, anchor=north west] (p0) at (0,0)
    {\color{white}Parte 0, Fundamentos comuns, obrigatória para todos};
  \foreach \c/\nome/\pt/\b/\i/\a in {
      0/Hardware/Parte 1/{Básico \cod{1.1} a \cod{1.5}}/{Intermediário \cod{1.6} a \cod{1.12}}/{Avançado \cod{1.13} a \cod{1.19}},
      1/Firmware/Parte 2/{Básico \cod{2.1} a \cod{2.6}}/{Intermediário \cod{2.7} a \cod{2.14}}/{Avançado \cod{2.15} a \cod{2.29}},
      2/Software/Parte 3/{Básico \cod{3.1} a \cod{3.4}}/{Intermediário \cod{3.5} a \cod{3.13}}/{Avançado \cod{3.14} a \cod{3.24}}}{
    \pgfmathsetmacro\x{57*\c}
    \node[nro moldura, minimum width=52mm, minimum height=72mm, anchor=north west] (m\c) at (\x,22) {};
    \node[nro rotulo, anchor=north] at (\x+26,26) {\rotulo{\pt}};
    \node[font=\fonttitulo\tamanho{13}\color{nro-tinta}, anchor=north] at (\x+26,31) {\nome};
    \node[nro claro, minimum width=42mm, minimum height=10mm] (b\c) at (\x+26,47) {\b};
    \node[nro medio, minimum width=42mm, minimum height=10mm] (i\c) at (\x+26,62) {\i};
    \node[nro forte, minimum width=42mm, minimum height=10mm] (a\c) at (\x+26,77) {\a};
    \draw[nro seta] (b\c.south) -- (i\c.north);
    \draw[nro seta] (i\c.south) -- (a\c.north);
    \draw[nro seta] (\x+26,11) -- (\x+26,22);
    \draw[nro seta] (\x+26,94) -- (\x+26,105);
  }
  \draw[nro ponte] (i0.east) -- (i1.west);
  \draw[nro ponte] (a0.east) -- (a1.west);
  \draw[nro ponte] (a1.east) -- (a2.west);
  \node[nro bloco, minimum width=166mm, minimum height=11mm, align=center, anchor=north west] at (0,105)
    {\color{white}Parte 4, Integração do sistema, trilha final comum};
\end{tikzpicture}

\Needspace{12\baselineskip}
\secao{Pontes entre trilhas}
Para atravessar para outra área, estude os tópicos da mesma linha juntos:

\par\vspace{1mm}
\begin{tabularx}{\linewidth}{@{}L{54mm}L{54mm}Y@{}}\reguatopo
  \cab{Hardware} & \cab{Firmware} & \cab{Software} \\ \reguameio
  \cod{1.7} a \cod{1.12}, fonte de corrente, sense, lógica, DAC, chaves e relés & \cod{2.9} a \cod{2.12}, periféricos, protocolos seriais e drivers & \vazio \\ \reguafina
  \cod{1.17}, proteções de hardware & \cod{2.28}, firmware de segurança crítica & \vazio \\ \reguafina
  \vazio & \cod{2.25}, BLE periférico & \cod{3.14} a \cod{3.16}, BLE central, bytes e contrato \\ \reguafina
  \vazio & \cod{2.22}, protocolo binário & \cod{3.15}, conversão de bytes \\ \reguabase
\end{tabularx}

%% ------------------------------------------------------------
\pagina{Estudo de caso}{O Connect em uma página}
O Connect é a plataforma de eletroestimulação da equipe e o laboratório desta apostila. Ele tem os blocos de qualquer dispositivo médico conectado: um aplicativo, um enlace sem fio, processadores com firmware e um estágio de potência que chega ao paciente. Cada bloco vive em um repositório e é estudado em uma trilha.

\par\vspace{3mm}\noindent
\begin{tikzpicture}[x=1mm, y=-1mm, node distance=0mm]
  \def\bw{78}\def\bh{14}\def\gap{17}
  \foreach \i/\t/\d/\r/\p in {
      0/App/{Flutter, celular}/{ErgoFES\_Connect\_software}/{Parte 3},
      1/ESP32/{Ponte BLE e UART, FreeRTOS}/{ergo\_fes\_connect\_ble\_firmware}/{Parte 2},
      2/ATSAMC20/{Gera o estímulo, C++20 bare-metal}/{ergo\_fes\_connect\_stimulus\_firmware}/{Parte 2},
      3/{Estágio de saída}/{DAC, fonte de corrente, +125 V, proteções}/{SCH070105 e SCH070102}/{Parte 1}}{
    \pgfmathsetmacro\y{\i*(\bh+\gap)}
    \node[nro bloco, minimum width=\bw mm, minimum height=\bh mm, anchor=north west, inner sep=3mm] (n\i) at (4,\y)
      {\fonttitulo\tamanho{11}\t\\[-0.4mm]{\color{nro-cortex-medio}\tamanho{8}\d}};
    \node[nro moldura, minimum width=68mm, minimum height=\bh mm, anchor=north west] (r\i) at (96,\y) {};
    \node[anchor=north west, inner sep=0] at (99,\y+2.2) {\color{nro-axon}\rotulo{\p}};
    \node[anchor=north west, inner sep=0] at (99,\y+7.4) {\ttfamily\tamanho{7.2}\color{nro-tinta}\r};
    \draw[nro-pia, line width=0.5pt] (n\i.east) -- (r\i.west);
  }
  \node[fill=nro-cortex-claro, minimum width=\bw mm, minimum height=\bh mm, anchor=north west, inner sep=3mm, align=left] (n4) at (4,{4*(\bh+\gap)})
    {\fonttitulo\tamanho{11}\color{nro-tinta}Paciente\\[-0.4mm]{\color{nro-cinza}\tamanho{8}Eletrodos}};
  \foreach \i/\j/\t in {0/1/Bluetooth Low Energy, 1/2/UART e WeProtocol, 2/3/SPI e sinais de controle, 3/4/Corrente de estímulo}{
    \draw[nro seta] (n\i.south west |- n\i.south) ++(8,0) coordinate (s) -- ($(n\j.north west)+(8,0)$)
      node[midway, right=2mm, nro rotulo] {\t};
  }
\end{tikzpicture}

\Needspace{12\baselineskip}
\secao{Repositórios e documentos de referência}
\par\vspace{1mm}
\begin{tabularx}{\linewidth}{@{}L{26mm}L{52mm}Y L{16mm}@{}}\reguatopo
  \cab{Área} & \cab{Referência} & \cab{Tecnologia} & \cab{Trilha} \\ \reguameio
  \textbf{Firmware stimulus} & \code{ergo\_fes\_connect\_stimulus\_firmware} & ATSAMC20N17A, C++20 bare-metal, MPLAB Harmony & Parte 2 \\ \reguafina
  \textbf{Firmware BLE} & \code{ergo\_fes\_connect\_ble\_firmware} & ESP32, ESP-IDF, Bluedroid, FreeRTOS & Parte 2 \\ \reguafina
  \textbf{Aplicativo} & \code{ErgoFES\_Connect\_software} (pasta com os 5 repositórios de software) & Flutter (\code{app\_main}), pacotes \code{bluetooth}, \code{byte\_conversion}, \code{date\_time\_helper}, \code{exception} & Parte 3 \\ \reguafina
  \textbf{Hardware} & \cod{SCH070105} (CFCU) e \cod{SCH070102} (CTCE) & Esquemáticos das placas & Parte 1 \\ \reguabase
\end{tabularx}

%% ------------------------------------------------------------
%% AS PARTES
%% ------------------------------------------------------------
% Parte 0: Fundamentos comuns
\begin{parte}{Fundamentos comuns}{A base que todo membro precisa, de qualquer área: como funciona um dispositivo de neuroestimulação, o mínimo de fisiologia e de segurança elétrica, e as ferramentas do dia a dia de engenharia.}
  \item Explicar a arquitetura de um eletroestimulador e o papel de cada camada.
  \item Trabalhar com binário, hexadecimal, terminal e Git sem travar.
  \item Ler datasheets, esquemáticos e normas com objetivo claro.
\end{parte}

\topico{Dispositivos de neuroestimulação}
Eletroestimuladores aplicam corrente controlada ao corpo para ativar nervos e músculos (FES, TENS) ou para modular a atividade do cérebro e do sistema vestibular (tDCS, GVS). Por trás de qualquer um deles há a mesma arquitetura: uma interface com o usuário, um enlace de comunicação, um microcontrolador que gera a forma de onda e um estágio de saída que entrega a corrente com segurança.

\begin{noproduto}
O Connect segue essa cadeia, nos dois sentidos: app, BLE, ESP32, UART, ATSAM, estágio de saída e paciente. O app vive em \code{ErgoFES\_Connect\_software}, a ponte BLE em \code{ergo\_fes\_connect\_ble\_firmware}, a geração do estímulo em \code{ergo\_fes\_connect\_stimulus\_firmware} e o estágio de saída nos esquemáticos \code{SCH070105} (CFCU) e \code{SCH070102} (CTCE).
\end{noproduto}

\begin{paraestudar}
  \item \leitura{https://en.wikipedia.org/wiki/Functional_electrical_stimulation}{Wikipedia: Functional electrical stimulation}
  \item \leitura{https://en.wikipedia.org/wiki/Transcranial_direct-current_stimulation}{Wikipedia: Transcranial direct-current stimulation}
  \item \leitura{https://en.wikipedia.org/wiki/Galvanic_vestibular_stimulation}{Wikipedia: Galvanic vestibular stimulation}
  \item Documento técnico interno do Connect e READMEs dos repositórios
\end{paraestudar}

\begin{pratique}
  \item Desenhe o diagrama de blocos de um eletroestimulador genérico e, ao lado, o do Connect. Indique em cada seta o protocolo ou sinal usado.
  \item Para cada repositório do Connect, escreva uma frase dizendo o que ele faz e com quem conversa.
  \item Explique com suas palavras a diferença entre um pulso de FES e a corrente contínua de tDCS.
\end{pratique}

\topico{Fisiologia básica para engenharia}
Nervo e músculo, potencial de ação, contração evocada por estímulo elétrico, a interface pele-eletrodo (e sua impedância) e o sistema vestibular. É o mínimo de biologia para entender por que os parâmetros de estímulo têm os limites que têm.

\begin{noproduto}
Esses conceitos explicam os parâmetros que o firmware permite configurar, como amplitude e largura de pulso, e por que o hardware precisa detectar eletrodo solto.
\end{noproduto}

\begin{paraestudar}
  \item \leitura{https://nba.uth.tmc.edu/neuroscience/}{Neuroscience Online (UTHealth), capítulos de neurônio e músculo}
  \item Livro: Princípios de Neurociências (Kandel et al.), capítulos introdutórios
\end{paraestudar}

\begin{pratique}
  \item Explique por que um pulso precisa de amplitude e duração mínimas para recrutar uma fibra (curva intensidade-duração).
  \item Modele pele e eletrodo como um circuito RC simples e esboce a tensão sobre ele sob um degrau de corrente.
\end{pratique}

\topico{Segurança elétrica e normas de equipamento médico}
Efeitos da corrente no corpo, as normas IEC 60601-1 (segurança geral de equipamento médico) e IEC 60601-2-10 (estimuladores de nervo e músculo), a gestão de risco pela ISO 14971 e o ciclo de vida de software médico pela IEC 62304.

\begin{noproduto}
O estágio de saída trabalha com +125 V. Por isso o Connect combina proteções em hardware (sobrecorrente, parada de emergência, eletrodo solto) com regras de firmware, como só habilitar a alta tensão em build Release.
\end{noproduto}

\begin{paraestudar}
  \item Normas IEC 60601-1, IEC 60601-2-10, ISO 14971 e IEC 62304 (verifique o acesso pela UFMG ou pela equipe)
  \item \leitura{https://www.allaboutcircuits.com/textbook/direct-current/chpt-3/}{All About Circuits: Electrical Safety}
\end{paraestudar}

\begin{pratique}
  \item Liste três perigos de um eletroestimulador e, para cada um, uma mitigação em hardware, em firmware ou no app.
  \item Encontre na IEC 60601-2-10 os limites aplicáveis à saída de um estimulador e compare com os do Connect.
\end{pratique}

\topico{Sistemas numéricos e representação de dados}
Binário, hexadecimal, complemento de dois, endianness e operações bit a bit. É a linguagem comum entre quem escreve firmware (registradores, protocolos) e quem escreve software (conversão de bytes).

\begin{noproduto}
Firmware e app trocam structs em binário. O \code{WeProtocol} no firmware e o pacote \code{byte\_conversion} no app precisam concordar byte a byte.
\end{noproduto}

\begin{paraestudar}
  \item \leitura{https://learn.sparkfun.com/tutorials/binary}{SparkFun: Binary}
  \item \leitura{https://learn.sparkfun.com/tutorials/hexadecimal}{SparkFun: Hexadecimal}
  \item \leitura{https://www.youtube.com/@BenEater}{Ben Eater no YouTube (vídeos sobre complemento de dois)}
\end{paraestudar}

\begin{pratique}
  \item Converta −1000 para int16 em complemento de dois e escreva os bytes em little-endian e em big-endian.
  \item Dado um registrador de 8 bits, escreva as máscaras para ligar o bit 3, desligar o bit 5 e testar o bit 7.
\end{pratique}

\topico{Linha de comando Linux}
Shell, navegação, permissões, variáveis de ambiente e PATH. Quase toda ferramenta de engenharia (compiladores, build systems, SDKs) é usada a partir do terminal.

\begin{noproduto}
CMake e Ninja no firmware stimulus, \code{idf.py} no firmware BLE e o Flutter SDK com FVM no app: as três bases de código são compiladas pelo terminal.
\end{noproduto}

\begin{paraestudar}
  \item \leitura{https://missing.csail.mit.edu/}{The Missing Semester of Your CS Education (MIT)}
  \item \leitura{https://linuxjourney.com/}{Linux Journey}
\end{paraestudar}

\begin{pratique}
  \item Adicione uma pasta ao PATH de forma permanente e explique em qual arquivo a configuração ficou.
  \item Escreva um script que compile um projeto e copie o binário gerado para uma pasta com a data no nome.
\end{pratique}

\topico{Git e fluxo de trabalho}
Branches, merge e rebase, submódulos, mensagens no padrão Conventional Commits e code review. É como uma equipe trabalha no mesmo código sem pisar no trabalho do outro.

\begin{noproduto}
O firmware usa a pasta \code{Modules/} como submódulo compartilhado e o app depende de pacotes locais (\code{we\_bluetooth}, \code{we\_exception}...). Saber mover um submódulo de commit é rotina.
\end{noproduto}

\begin{paraestudar}
  \item \leitura{https://git-scm.com/book/pt-br/v2}{Pro Git (em português), capítulos 1 a 3 e a seção de submódulos}
  \item \leitura{https://learngitbranching.js.org/?locale=pt_BR}{Learn Git Branching}
  \item \leitura{https://www.conventionalcommits.org/pt-br/v1.0.0/}{Conventional Commits}
\end{paraestudar}

\begin{pratique}
  \item Clone um repositório com submódulos, mude o submódulo de commit e explique o que aparece no \code{git} \code{status} do repositório pai.
  \item Faça um rebase interativo juntando três commits em um, com mensagem no padrão Conventional Commits.
  \item Revise o PR de um colega e deixe ao menos dois comentários construtivos.
\end{pratique}

\topico{Leitura de documentação técnica}
Como ler datasheet, reference manual, esquemático e mapa de registradores: onde estão os limites absolutos, as condições típicas, os diagramas de tempo e os valores de reset. Ler com uma pergunta em mente economiza horas.

\begin{noproduto}
Os esquemáticos do Connect e os datasheets do ATSAMC20, do DAC53204 e do BQ25792 são as leituras de referência das trilhas de hardware e firmware.
\end{noproduto}

\begin{paraestudar}
  \item \leitura{https://learn.sparkfun.com/tutorials/how-to-read-a-schematic}{SparkFun: How to Read a Schematic}
  \item Datasheet do DAC53204 (Texas Instruments) como caso de estudo
\end{paraestudar}

\begin{pratique}
  \item No datasheet do DAC53204, encontre a tensão máxima de alimentação, o tempo de acomodação e o formato do quadro SPI.
  \item Escolha um registrador do ATSAMC20 e explique cada campo e o seu valor de reset.
\end{pratique}

% Parte 1: Hardware (gerado a partir da apostila NeuroDev; edite à vontade)
\begin{parte}{Hardware}{Da lei de Ohm ao estágio de saída de alta tensão. A trilha ensina eletrônica analógica, digital e de potência de forma geral e usa os esquemáticos do Connect como material de estudo real.}
  \item Analisar e dimensionar circuitos analógicos e digitais de uma placa real.
  \item Entender fontes chaveadas, gate drivers, bateria e proteções.
  \item Planejar e executar a validação de um hardware em bancada.
\end{parte}

\nivel{Básico}{Circuitos, componentes, bancada e leitura de esquemático.}

\topico{Circuitos elétricos}
Lei de Ohm, leis de Kirchhoff, divisor de tensão, potência e dissipação. É a base para dimensionar qualquer componente de uma placa.

\begin{noproduto}
Divisores de tensão, resistores de sense e a dissipação dos componentes do estágio de saída são dimensionados com essas ferramentas.
\end{noproduto}

\begin{paraestudar}
  \item \leitura{https://learn.sparkfun.com/tutorials/voltage-current-resistance-and-ohms-law}{SparkFun: Voltage, Current, Resistance, and Ohm's Law}
  \item \leitura{https://learn.sparkfun.com/tutorials/voltage-dividers}{SparkFun: Voltage Dividers}
  \item \leitura{https://www.allaboutcircuits.com/textbook/}{All About Circuits: textbook, volume de corrente contínua}
\end{paraestudar}

\begin{pratique}
  \item Encontre um divisor de tensão no esquemático e calcule a tensão de saída e a corrente quiescente.
  \item Calcule a potência dissipada no resistor de sense da fonte de corrente na corrente máxima de estímulo.
\end{pratique}

\topico{Componentes passivos}
Resistor (potência e tolerância), capacitor (tipos, ESR, desacoplamento, perda de capacitância com tensão DC) e indutor (saturação e corrente RMS). Escolher o passivo certo é metade de um projeto confiável.

\begin{noproduto}
Cada capacitor dos esquemáticos tem um motivo para seu tipo e sua tensão nominal, principalmente perto dos trilhos de alta tensão.
\end{noproduto}

\begin{paraestudar}
  \item \leitura{https://learn.sparkfun.com/tutorials/capacitors}{SparkFun: Capacitors}
  \item Livro: The Art of Electronics (Horowitz e Hill), capítulo 1
\end{paraestudar}

\begin{pratique}
  \item Escolha três capacitores do esquemático e justifique o tipo (cerâmico, eletrolítico) e a tensão nominal de cada um.
  \item Explique por que um cerâmico X7R de 10 µF perde capacitância quando recebe tensão DC.
\end{pratique}

\topico{Semicondutores}
Diodo, zener, LED, BJT (inclusive o par Darlington) e MOSFET como chave: regiões de operação, ganho, tensão de limiar e perdas.

\begin{noproduto}
O estágio de saída usa transistores como o Darlington FJD5304D. Entender saturação e ganho mínimo é o que garante que eles conduzam a corrente prevista.
\end{noproduto}

\begin{paraestudar}
  \item \leitura{https://learn.sparkfun.com/tutorials/diodes}{SparkFun: Diodes}
  \item \leitura{https://learn.sparkfun.com/tutorials/transistors}{SparkFun: Transistors}
  \item Datasheet do FJD5304D
\end{paraestudar}

\begin{pratique}
  \item Calcule a corrente de base necessária para saturar um BJT do esquemático a partir do ganho mínimo do datasheet.
  \item Compare o uso de um MOSFET e de um Darlington como chave em alta tensão.
\end{pratique}

\topico{Instrumentação de bancada}
Multímetro, osciloscópio (trigger, acoplamento, largura de banda), fonte de bancada com limite de corrente, carga resistiva e ponteiras de alta tensão e diferenciais. Medir errado é pior que não medir.

\begin{noproduto}
Medir pulsos no estágio de saída exige carga resistiva e cuidado com a referência: GND e GND\_HV não são o mesmo ponto.
\end{noproduto}

\begin{paraestudar}
  \item \leitura{https://learn.sparkfun.com/tutorials/how-to-use-a-multimeter}{SparkFun: How to Use a Multimeter}
  \item \leitura{https://learn.sparkfun.com/tutorials/how-to-use-an-oscilloscope}{SparkFun: How to Use an Oscilloscope}
  \item \leitura{https://www.youtube.com/@EEVblog}{EEVblog no YouTube}
\end{paraestudar}

\begin{pratique}
  \item Meça um pulso bifásico numa carga resistiva e registre amplitude, largura e intervalo.
  \item Explique por que não se deve prender a referência da ponteira no GND\_HV sem isolação ou ponteira diferencial.
\end{pratique}

\topico{Leitura de esquemático}
Hierarquia de páginas, net labels, designators, notas de montagem e grounds separados. Saber navegar num esquemático grande é a habilidade mais usada no dia a dia.

\begin{noproduto}
Os esquemáticos \code{SCH070105} (CFCU) e \code{SCH070102} (CTCE) têm várias páginas e dois grounds, GND e GND\_HV.
\end{noproduto}

\begin{paraestudar}
  \item \leitura{https://learn.sparkfun.com/tutorials/how-to-read-a-schematic}{SparkFun: How to Read a Schematic}
  \item Esquemáticos SCH070105 e SCH070102
\end{paraestudar}

\begin{pratique}
  \item Siga um sinal de controle desde o pino do ATSAM até o estágio de saída, listando cada página por onde ele passa.
  \item Liste onde GND e GND\_HV se encontram (ou não) e explique o motivo.
\end{pratique}

\nivel{Intermediário}{Op-amps, fontes de corrente, conversores e chaves.}

\topico{Amplificadores operacionais}
Op-amp ideal e real (offset, GBW, slew rate, saída rail-to-rail), buffer, inversor e não inversor, comparador, histerese e Schmitt trigger.

\begin{noproduto}
Op-amps aparecem como fonte de corrente, como comparadores das proteções e como buffers de sinal ao longo das placas.
\end{noproduto}

\begin{paraestudar}
  \item Livro: Op Amps for Everyone (Texas Instruments, Ron Mancini), gratuito
  \item \leitura{https://www.allaboutcircuits.com/textbook/}{All About Circuits: textbook, capítulo de op-amps}
\end{paraestudar}

\begin{pratique}
  \item Projete um comparador com histerese de 100 mV e simule no LTspice.
  \item Encontre no esquemático um op-amp em cada configuração (buffer, não inversor, comparador) e calcule o ganho ou os limiares.
\end{pratique}

\topico{Fontes de corrente controladas por tensão}
Um op-amp, um transistor e um resistor de sense formam uma malha que força a tensão no shunt a seguir a referência. O resultado é I = Vref / Rsense, independente da carga até o limite de compliance. É o coração de qualquer estimulador por corrente.

\begin{noproduto}
No Connect, a fonte de corrente é o conjunto TLV9151 + Q36 + R196, com a referência vinda do DAC.
\end{noproduto}

\begin{paraestudar}
  \item Livro: The Art of Electronics, seção de fontes de corrente com op-amp
  \item Datasheet do TLV9151
\end{paraestudar}

\begin{pratique}
  \item Calcule a corrente de saída para Vref = 1 V com o valor de R196 do esquemático.
  \item Encontre a tensão de compliance: até que impedância de carga a fonte mantém a corrente com +125 V?
  \item Simule o circuito no LTspice com uma carga RC representando pele e eletrodo.
\end{pratique}

\topico{Medição de corrente}
Resistor shunt, amplificador de sense de corrente e comparador com latch para desligamento por sobrecorrente. Medir corrente com precisão é o que permite controlá-la e limitá-la.

\begin{noproduto}
O INA310 combina amplificador de sense e comparador, e alimenta a proteção de sobrecorrente.
\end{noproduto}

\begin{paraestudar}
  \item Datasheet do INA310 (Texas Instruments)
  \item TI Precision Labs, série de vídeos sobre current sensing (site da TI)
\end{paraestudar}

\begin{pratique}
  \item Calcule a tensão de saída do amplificador de sense na corrente máxima e verifique se cabe na faixa do ADC.
  \item Explique por que o comparador de proteção precisa travar (latch) em vez de apenas comparar.
\end{pratique}

\topico{Lógica digital}
Portas lógicas, buffers e inversores Schmitt, saídas open-drain com pull-up e compatibilidade de níveis lógicos entre 3,3 V e 5 V.

\begin{noproduto}
Buffers e inversores como o 74LVC1G14 e o 74LVC1G17 limpam sinais que cruzam a placa e montam a lógica de desligamento.
\end{noproduto}

\begin{paraestudar}
  \item \leitura{https://learn.sparkfun.com/tutorials/logic-levels}{SparkFun: Logic Levels}
  \item \leitura{https://learn.sparkfun.com/tutorials/pull-up-resistors}{SparkFun: Pull-up Resistors}
  \item Datasheets 74LVC1G14 e 74LVC1G17
\end{paraestudar}

\begin{pratique}
  \item Dimensione o pull-up de uma linha open-drain para um tempo de subida de 1 µs com 50 pF.
  \item Explique por que usar um buffer Schmitt num sinal que vem de outra placa.
\end{pratique}

\topico{Conversores DAC e ADC}
Resolução, LSB, tensão de referência (VDD, interna ou externa), erros de offset e ganho, INL e DNL. Conversores são a fronteira entre o mundo digital do firmware e o analógico do circuito.

\begin{noproduto}
O DAC53204 define a referência da fonte de corrente, ou seja, a amplitude do estímulo.
\end{noproduto}

\begin{paraestudar}
  \item \leitura{https://learn.sparkfun.com/tutorials/analog-to-digital-conversion}{SparkFun: Analog to Digital Conversion}
  \item Datasheet do DAC53204 (Texas Instruments)
\end{paraestudar}

\begin{pratique}
  \item Calcule o passo de corrente na saída correspondente a 1 LSB do DAC.
  \item Compare a precisão da saída usando referência VDD e referência interna.
\end{pratique}

\topico{Chaves analógicas e multiplexação}
Resistência on, carga injetada e faixa de sinal das chaves analógicas, usadas para rotear um sinal entre caminhos diferentes.

\begin{noproduto}
O NLAS4157 compartilha o sinal entre os modos FES e tDCS.
\end{noproduto}

\begin{paraestudar}
  \item Datasheet do NLAS4157 (onsemi)
\end{paraestudar}

\begin{pratique}
  \item Encontre o NLAS4157 no esquemático e descreva quais sinais ele seleciona e quem o controla.
  \item Estime o erro introduzido pela resistência on da chave no caminho do sinal.
\end{pratique}

\topico{Relés de estado sólido e isolação}
Relés PhotoMOS: LED de entrada, MOSFETs de saída, isolação galvânica e acionamento por transistor. Permitem abrir um caminho de alta tensão a partir de um sinal lógico.

\begin{noproduto}
O AQW280SX é acionado via transistor e abre o caminho de saída quando necessário.
\end{noproduto}

\begin{paraestudar}
  \item Datasheet do AQW280SX (Panasonic)
\end{paraestudar}

\begin{pratique}
  \item Calcule o resistor de entrada do relé para a corrente de LED recomendada.
  \item Explique qual falha o relé previne quando aberto e o que acontece se ele colar fechado.
\end{pratique}

\nivel{Avançado}{Potência, bateria, proteções, layout e validação.}

\topico{Fontes chaveadas}
Topologias buck, boost e push-pull, malha de controle, transformador e retificação. É como se gera, com eficiência, uma tensão muito maior ou menor que a da bateria.

\begin{noproduto}
Os +125 V do estágio de saída vêm de um conversor push-pull com o controlador LM5030.
\end{noproduto}

\begin{paraestudar}
  \item Livro: Fundamentals of Power Electronics (Erickson e Maksimović), capítulos 1 a 6
  \item Nota de aplicação TI SLVA477: Basic Calculation of a Buck Converter's Power Stage
  \item Datasheet do LM5030
\end{paraestudar}

\begin{pratique}
  \item Desenhe o caminho de energia da bateria até os +125 V, indicando cada conversor.
  \item Calcule o ciclo de trabalho e a relação de espiras esperados no push-pull.
\end{pratique}

\topico{Gate drivers e ponte H}
Half-bridge e ponte H, gate drivers com bootstrap e seus limites: o capacitor de bootstrap precisa ser recarregado, então não sustenta condução contínua indefinidamente. A ponte H permite inverter a polaridade da corrente na carga.

\begin{noproduto}
O IRS25752 aciona o estágio bifásico do FES. O tDCS, por ser corrente contínua, exige outra solução.
\end{noproduto}

\begin{paraestudar}
  \item Infineon (IR) AN-978: HV Floating MOS-Gate Driver ICs
  \item Datasheet do IRS25752
\end{paraestudar}

\begin{pratique}
  \item Calcule o capacitor de bootstrap mínimo para a maior largura de pulso usada.
  \item Explique por que corrente contínua de longa duração não combina com bootstrap.
\end{pratique}

\topico{Gerenciamento de bateria}
Química Li-ion (carga CC/CV, limites de temperatura), carregadores controlados por I²C, power path e proteções. Bateria mal gerenciada é risco de segurança, não só de autonomia.

\begin{noproduto}
O carregador é o BQ25792, configurado pelo firmware via I²C.
\end{noproduto}

\begin{paraestudar}
  \item \leitura{https://batteryuniversity.com/}{Battery University}
  \item Datasheet do BQ25792 (Texas Instruments)
\end{paraestudar}

\begin{pratique}
  \item Liste os registradores do BQ25792 que definem corrente de carga, tensão final e limite de entrada.
  \item Explique o que o power path faz quando o carregador está conectado e a bateria está descarregada.
\end{pratique}

\topico{Sequenciamento de energia e load switches}
Habilitação ordenada dos trilhos de alimentação, corrente de inrush, soft-start e descarga ativa ao desligar. A ordem em que a placa liga define se ela liga direito.

\begin{noproduto}
Os trilhos do Connect, inclusive o de alta tensão, são habilitados em sequência.
\end{noproduto}

\begin{paraestudar}
  \item Notas de aplicação da TI sobre load switches (site da TI)
\end{paraestudar}

\begin{pratique}
  \item Monte a tabela de sequência de ligação dos trilhos a partir do esquemático.
  \item Meça com o osciloscópio a subida de dois trilhos e compare com a ordem esperada.
\end{pratique}

\topico{Proteções de hardware}
Sobrecorrente, parada de emergência, detecção de eletrodo solto e uma lógica OR que desliga a saída sem depender do firmware. Em equipamento médico, a última linha de defesa precisa ser física.

\begin{noproduto}
As proteções do Connect se juntam numa lógica OR de desligamento, documentada nas páginas de proteção do \code{SCH070102}.
\end{noproduto}

\begin{paraestudar}
  \item IEC 60601-2-10, requisitos de proteção de estimuladores
  \item Esquemático SCH070102, páginas de proteção
\end{paraestudar}

\begin{pratique}
  \item Desenhe a lógica OR de desligamento: quais entradas, qual saída e o que ela comanda.
  \item Para cada proteção, indique o tempo de resposta esperado e quem rearma o sistema.
\end{pratique}

\topico{Integridade de sinal e layout}
Desacoplamento, laços de corrente, planos de referência, distâncias de isolação (creepage e clearance) e EMC. Um esquemático certo pode falhar por causa de um layout ruim.

\begin{noproduto}
A placa convive com +125 V e lógica de 3,3 V, então as distâncias de isolação seguem a IEC 60601-1.
\end{noproduto}

\begin{paraestudar}
  \item \leitura{https://www.youtube.com/@PhilsLab}{Phil's Lab no YouTube}
  \item Palestras de Rick Hartley sobre aterramento e retorno de corrente (YouTube)
  \item Livro: Electromagnetic Compatibility Engineering (Henry Ott)
\end{paraestudar}

\begin{pratique}
  \item No layout, encontre o laço de maior di/dt do conversor e proponha como reduzi-lo.
  \item Verifique a distância de isolação entre a região de alta tensão e a de baixa tensão e compare com a tabela da IEC 60601-1.
\end{pratique}

\topico{Validação de hardware}
Plano de teste, teste de disparo de cada proteção e caracterização da saída (linearidade, erro, ripple). Validar é provar, com dados, que o hardware faz o que o projeto promete.

\begin{noproduto}
A caracterização principal é a curva corrente de saída × código do DAC, medida em carga resistiva.
\end{noproduto}

\begin{paraestudar}
  \item Modelo de plano de teste da equipe (a criar)
\end{paraestudar}

\begin{pratique}
  \item Escreva um plano de teste para a proteção de sobrecorrente: setup, passos e critério de aprovação.
  \item Levante a curva corrente de saída × código do DAC e calcule o erro máximo.
\end{pratique}

% Parte 2: Firmware (gerado a partir da apostila NeuroDev; edite à vontade)
\begin{parte}{Firmware}{Programação de microcontroladores do registrador ao sistema completo: C e C++ embarcado, periféricos, protocolos, arquitetura, tempo real, Bluetooth e segurança. Os dois processadores do Connect, um ATSAMC20 bare-metal e um ESP32 com FreeRTOS, servem de laboratório.}
  \item Escrever código embarcado seguro, sem heap e sem bloqueio, em C e C++.
  \item Configurar periféricos, escrever drivers e depurar direto no hardware.
  \item Projetar arquitetura, protocolos e testes de um firmware crítico.
\end{parte}

\nivel{Básico}{C, C++ embarcado, microcontrolador, GPIO e ferramentas de build.}

\topico{Linguagem C}
Tipos de largura fixa (\code{uint8\_t}, \code{int16\_t}), ponteiros, structs, \code{volatile}, \code{static}, bitfields e máscaras. É a base de qualquer código que toca registradores.

\begin{noproduto}
Mesmo com o firmware escrito em C++20, a interface com o hardware (registradores e código gerado pelo Harmony) é C puro.
\end{noproduto}

\begin{paraestudar}
  \item \leitura{https://beej.us/guide/bgc/}{Beej's Guide to C Programming}
  \item Livro: The C Programming Language (Kernighan e Ritchie)
\end{paraestudar}

\begin{pratique}
  \item Escreva uma função que empacote três valores (4, 6 e 6 bits) num \code{uint16\_t} e outra que os desempacote.
  \item Mostre, com um exemplo, o que pode acontecer se uma flag alterada numa interrupção não for \code{volatile}.
\end{pratique}

\topico{C++ para embarcados}
Classes, namespaces, \code{enum class}, \code{constexpr}, templates, \code{std::array} e \code{std::span} e RAII, usados sem exceções, sem RTTI e sem alocação dinâmica. C++ bem usado dá segurança de tipos sem custo de desempenho.

\begin{noproduto}
O \code{ergo\_fes\_connect\_stimulus\_firmware} é C++20 bare-metal, sem exceções, sem RTTI e sem heap.
\end{noproduto}

\begin{paraestudar}
  \item \leitura{https://en.cppreference.com/}{cppreference}
  \item \leitura{https://embeddedartistry.com/}{Embedded Artistry, artigos sobre C++ em embarcados}
\end{paraestudar}

\begin{pratique}
  \item Reescreva um conjunto de \code{\#define} de estados como \code{enum class} e mostre o erro de compilação que isso passa a evitar.
  \item Implemente um guard RAII que desabilita interrupções no construtor e as restaura no destrutor.
\end{pratique}

\topico{Arquitetura de microcontroladores}
CPU, memória flash e RAM, barramento, periféricos mapeados em memória e árvore de clock. Entender o mapa do chip é entender o que o código realmente faz.

\begin{noproduto}
O processador de estímulo é o ATSAMC20N17A, um Cortex-M0+ da Microchip.
\end{noproduto}

\begin{paraestudar}
  \item \leitura{https://www.youtube.com/@StateMachineCOM}{Quantum Leaps (Miro Samek): Modern Embedded Systems Programming}
  \item Datasheet do SAM C20/C21 (Microchip), visão geral e mapa de memória
\end{paraestudar}

\begin{pratique}
  \item A partir do datasheet, desenhe o mapa de memória do ATSAMC20N17A (flash, RAM, periféricos).
  \item Explique como o clock chega a um SERCOM: fonte, gerador de clock e habilitação.
\end{pratique}

\topico{GPIO}
Entrada, saída, pull-up interno, leitura de botões com debounce e acesso direto aos registradores de porta. O primeiro contato real com o hardware.

\begin{noproduto}
No ATSAMC20 os pinos são controlados pelo periférico PORT.
\end{noproduto}

\begin{paraestudar}
  \item \leitura{https://learn.sparkfun.com/tutorials/pull-up-resistors}{SparkFun: Pull-up Resistors}
  \item Datasheet do SAM C20, capítulo PORT
\end{paraestudar}

\begin{pratique}
  \item Pisque um LED escrevendo direto nos registradores do PORT, sem bibliotecas.
  \item Implemente um debounce não bloqueante, baseado em tempo, para um botão.
\end{pratique}

\topico{Toolchain}
Compilador, assembler e linker; \code{arm-none-eabi-gcc}, arquivos objeto, ELF, binário e map file. Saber o que cada etapa produz ajuda a entender erros de build e problemas de memória.

\begin{noproduto}
O firmware stimulus é compilado com \code{arm-none-eabi-gcc} e gera um ELF com map file.
\end{noproduto}

\begin{paraestudar}
  \item \leitura{https://interrupt.memfault.com/blog/zero-to-main-1}{Memfault Interrupt: From Zero to main()}
\end{paraestudar}

\begin{pratique}
  \item Compile o firmware e use \code{arm-none-eabi-size} para ver quanto de flash e RAM é usado.
  \item Abra o map file e encontre as três funções que mais ocupam flash.
\end{pratique}

\topico{Build system}
CMake, Ninja, toolchain files e \code{compile\_commands.json}, o arquivo que alimenta clangd e clang-tidy. Um build reprodutível é o que permite que qualquer pessoa compile o projeto.

\begin{noproduto}
O firmware stimulus usa CMake e Ninja com um toolchain file para ARM.
\end{noproduto}

\begin{paraestudar}
  \item \leitura{https://cmake.org/cmake/help/latest/guide/tutorial/index.html}{Tutorial oficial do CMake}
\end{paraestudar}

\begin{pratique}
  \item Configure e compile o firmware do zero numa pasta de build nova, nas configurações Debug e Release.
  \item Explique o que o toolchain file do projeto define e por que ele é necessário.
\end{pratique}

\nivel{Intermediário}{Cortex-M, periféricos, protocolos seriais, drivers e depuração.}

\topico{Cortex-M0+}
Registradores do núcleo, modos de execução, NVIC, interrupções e prioridades, exceções e SysTick. É o comportamento comum a todos os microcontroladores dessa família.

\begin{noproduto}
O ATSAMC20 é um Cortex-M0+; tudo deste tópico vale diretamente para ele.
\end{noproduto}

\begin{paraestudar}
  \item ARM: Cortex-M0+ Devices Generic User Guide
  \item Livro: The Definitive Guide to ARM Cortex-M0 and Cortex-M0+ Processors (Joseph Yiu)
\end{paraestudar}

\begin{pratique}
  \item Configure o SysTick para gerar uma interrupção a cada 1 ms e conte o tempo desde o boot.
  \item Explique o que acontece quando duas interrupções de mesma prioridade chegam ao mesmo tempo.
\end{pratique}

\topico{Startup e linker script}
Tabela de vetores, handler de reset, cópia de \code{.data}, zeragem de \code{.bss} e o linker script que diz onde cada coisa mora na memória. É o que acontece antes do \code{main()}.

\begin{noproduto}
O firmware stimulus tem seu próprio \code{Startup.c} e arquivo \code{.ld}.
\end{noproduto}

\begin{paraestudar}
  \item \leitura{https://interrupt.memfault.com/blog/zero-to-main-1}{Memfault Interrupt: série From Zero to main()}
\end{paraestudar}

\begin{pratique}
  \item Leia o \code{Startup.c} do firmware e descreva cada passo até chegar ao \code{main()}.
  \item Crie uma variável numa seção própria pelo linker script e confirme o endereço no map file.
\end{pratique}

\topico{Periféricos de um microcontrolador}
Portas, interfaces seriais configuráveis (SPI, I²C, USART), timers, ADC, EEPROM e controlador de memória não volátil. Cada periférico tem registradores de configuração, de status e de dados.

\begin{noproduto}
No ATSAMC20: PORT, SERCOM (SPI, I²C, USART), TC (timers), ADC, EEPROM e NVM.
\end{noproduto}

\begin{paraestudar}
  \item Datasheet do SAM C20/C21 (Microchip), capítulos de cada periférico
\end{paraestudar}

\begin{pratique}
  \item Liste quais SERCOMs o firmware usa, em qual modo e para qual CI externo.
  \item Configure um timer para gerar uma onda quadrada de 1 kHz num pino e meça no osciloscópio.
\end{pratique}

\topico{Geradores de código de fabricante}
Ferramentas como MPLAB Harmony/MCC geram o BSP e as bibliotecas de periféricos (PLIBs) a partir de uma configuração gráfica. A regra de ouro: código gerado não se edita à mão; muda-se a configuração e regenera.

\begin{noproduto}
O firmware stimulus usa MPLAB Harmony/MCC, com fronteira clara entre o código gerado e o código da equipe.
\end{noproduto}

\begin{paraestudar}
  \item \leitura{https://github.com/Microchip-MPLAB-Harmony}{MPLAB Harmony no GitHub (repositórios csp e documentação)}
\end{paraestudar}

\begin{pratique}
  \item Altere a frequência de um SPI pelo MCC, regenere e mostre o diff resultante.
  \item Explique onde fica a fronteira entre código gerado e código da equipe no repositório.
\end{pratique}

\topico{Protocolos seriais}
SPI, I²C e UART: modos de clock, endereçamento, framing e checksum. São as estradas por onde um microcontrolador conversa com CIs e com outros processadores.

\begin{noproduto}
SPI fala com o DAC, I²C com o carregador e a UART liga o ATSAM ao ESP32.
\end{noproduto}

\begin{paraestudar}
  \item \leitura{https://learn.sparkfun.com/tutorials/serial-peripheral-interface-spi}{SparkFun: Serial Peripheral Interface (SPI)}
  \item \leitura{https://learn.sparkfun.com/tutorials/i2c}{SparkFun: I2C}
  \item \leitura{https://learn.sparkfun.com/tutorials/serial-communication}{SparkFun: Serial Communication}
\end{paraestudar}

\begin{pratique}
  \item Capture com um analisador lógico uma escrita SPI no DAC e decodifique o quadro.
  \item Capture um quadro na UART entre ATSAM e ESP32 e identifique cabeçalho, payload e checksum.
\end{pratique}

\topico{Drivers de CIs externos}
Como escrever um driver a partir do mapa de registradores: camada de transporte, definição de registradores, API de alto nível e tratamento de erro. Um bom driver esconde o protocolo e expõe intenções.

\begin{noproduto}
Os drivers do DAC53204 e do BQ25792 são os modelos de referência da equipe.
\end{noproduto}

\begin{paraestudar}
  \item Datasheets do DAC53204 e do BQ25792
  \item Drivers existentes no repositório
\end{paraestudar}

\begin{pratique}
  \item Leia o driver do DAC53204 e desenhe a sequência de comandos de inicialização.
  \item Implemente a leitura de um registrador de status do BQ25792 que ainda não tenha função no driver.
\end{pratique}

\topico{Gravação e depuração}
SWD, J-Link, OpenOCD e GDB: breakpoints, watchpoints e inspeção de registradores e memória. Depurador bem usado substitui horas de \code{printf}.

\begin{noproduto}
O ATSAMC20 é gravado e depurado por SWD.
\end{noproduto}

\begin{paraestudar}
  \item \leitura{https://wiki.segger.com/}{SEGGER Wiki}
  \item \leitura{https://openocd.org/}{OpenOCD}
  \item \leitura{https://sourceware.org/gdb/documentation/}{Documentação do GDB}
\end{paraestudar}

\begin{pratique}
  \item Grave o firmware, pare num breakpoint dentro de um Handler e inspecione uma variável e um registrador de periférico.
  \item Use um watchpoint para descobrir quem altera um campo do \code{SystemData}.
\end{pratique}

\topico{Qualidade de código}
clang-format, clang-tidy, warnings tratados como erros, padrão de nomes e comentários com autor. Qualidade é o que torna o código revisável por outra pessoa.

\begin{noproduto}
O firmware usa \code{.clang-format}, clang-tidy e warnings como erros.
\end{noproduto}

\begin{paraestudar}
  \item \leitura{https://clang.llvm.org/docs/ClangFormat.html}{ClangFormat}
  \item \leitura{https://clang.llvm.org/extra/clang-tidy/}{clang-tidy}
  \item Guia de estilo interno do repositório
\end{paraestudar}

\begin{pratique}
  \item Rode clang-tidy no projeto e corrija dois avisos, explicando cada um.
  \item Configure o editor para formatar ao salvar usando o \code{.clang-format} do repositório.
\end{pratique}

\nivel{Avançado}{Arquitetura, tempo real, BLE, protocolo próprio, segurança e testes.}

\topico{Arquitetura superloop cooperativa}
Services com \code{Init}, \code{DeInit} e \code{Handler} chamados em sequência no laço principal. Funciona desde que nenhum Handler bloqueie: cada um tem um orçamento de latência.

\begin{noproduto}
O firmware stimulus é organizado assim; o \code{main} chama os Handlers de todos os services a cada volta.
\end{noproduto}

\begin{paraestudar}
  \item Livro: Making Embedded Systems (Elecia White), capítulos de arquitetura
  \item Código do main e dos services no firmware stimulus
\end{paraestudar}

\begin{pratique}
  \item Meça com um GPIO e o osciloscópio o tempo de uma volta do superloop e o do Handler mais lento.
  \item Converta um trecho com espera ativa (\code{delay}) numa versão não bloqueante.
\end{pratique}

\topico{Máquinas de estados}
Modelagem de estados e transições, implementação com \code{switch} ou tabela, e diagramas. É a forma mais clara de descrever comportamento que depende do histórico.

\begin{noproduto}
O \code{StimulusService} e o \code{SystemController} são máquinas de estados.
\end{noproduto}

\begin{paraestudar}
  \item Livro: Practical UML Statecharts in C/C++ (Miro Samek)
  \item \leitura{https://mermaid.js.org/}{Mermaid: diagramas de estado}
\end{paraestudar}

\begin{pratique}
  \item Desenhe o diagrama de estados do \code{StimulusService} a partir do código.
  \item Encontre uma transição não tratada, ou justifique por que todas estão.
\end{pratique}

\topico{Temporização precisa}
Base de tempo em microssegundos, geração de pulsos com largura e intervalo exatos e medição de jitter. Em estimulação, tempo é parâmetro clínico.

\begin{noproduto}
O \code{MicroTimerService} fornece a base de µs usada para gerar os pulsos bifásicos.
\end{noproduto}

\begin{paraestudar}
  \item Datasheet do SAM C20, capítulo TC
  \item Código do MicroTimerService
\end{paraestudar}

\begin{pratique}
  \item Meça o jitter da largura do pulso em 1000 pulsos e relacione com a carga do superloop.
  \item Explique como o firmware trata o overflow do contador de µs.
\end{pratique}

\topico{Concorrência com interrupções}
Seções críticas, condições de corrida, acesso atômico a variáveis multibyte e buffers circulares entre interrupção e laço principal. Bugs de concorrência são raros, intermitentes e caros.

\begin{noproduto}
Num superloop, dados que chegam por interrupção, como os bytes da UART, precisam de um caminho seguro até o laço principal.
\end{noproduto}

\begin{paraestudar}
  \item \leitura{https://embeddedartistry.com/}{Embedded Artistry: Creating a Circular Buffer in C and C++}
  \item \leitura{https://interrupt.memfault.com/}{Memfault Interrupt, artigos sobre interrupções no Cortex-M}
\end{paraestudar}

\begin{pratique}
  \item Implemente um buffer circular de um produtor e um consumidor sem desabilitar interrupções e justifique por que é seguro.
  \item Crie de propósito uma condição de corrida numa variável de 32 bits lida em duas partes e reproduza o erro.
\end{pratique}

\topico{Tratamento de erros}
Propagação de erro sem exceções, guards de escopo e transição para estados de falha seguros. Todo erro precisa de um destino definido.

\begin{noproduto}
O firmware usa o tipo \code{ReturnStatus} e guards de escopo.
\end{noproduto}

\begin{paraestudar}
  \item Código de ReturnStatus e dos guards no repositório
\end{paraestudar}

\begin{pratique}
  \item Rastreie o que acontece desde uma falha de escrita SPI no DAC até o estado do sistema.
  \item Proponha um código de erro novo e liste todos os lugares que precisariam tratá-lo.
\end{pratique}

\topico{Persistência}
EEPROM e memória não volátil, layout versionado dos dados, escrita atrasada para poupar ciclos e desgaste. Dados persistidos sobrevivem a versões de firmware, então precisam de versão.

\begin{noproduto}
O firmware stimulus persiste dados em EEPROM usando escrita atrasada.
\end{noproduto}

\begin{paraestudar}
  \item Datasheet do SAM C20, capítulos NVMCTRL e RWW EEPROM
\end{paraestudar}

\begin{pratique}
  \item Estime quantos anos a EEPROM dura com a política atual de escrita.
  \item Explique como o firmware deve reagir a dados gravados por uma versão antiga.
\end{pratique}

\topico{Estado central e sincronização}
Uma fonte única do estado do dispositivo e a forma de espelhá-lo em outros processadores. Evita que cada módulo tenha sua própria versão da verdade.

\begin{noproduto}
O singleton \code{SystemData} concentra o estado, e os valores são espelhados no ESP32 para leitura via BLE.
\end{noproduto}

\begin{paraestudar}
  \item Código do SystemData e da camada de sincronização
\end{paraestudar}

\begin{pratique}
  \item Liste os campos do \code{SystemData} e, para cada um, quem escreve e quem lê.
  \item Desenhe a sequência de uma mudança no nível da bateria até o valor chegar ao app.
\end{pratique}

\topico{Protocolo binário próprio}
Framing, serialização de structs, identificadores de mensagem, versionamento e regras de compatibilidade. Um protocolo vive mais que qualquer versão de firmware ou de app.

\begin{noproduto}
O \code{WeProtocol} liga ATSAM, ESP32 e app, com wire IDs que precisam permanecer compatíveis.
\end{noproduto}

\begin{paraestudar}
  \item Documentação e código do WeProtocol no repositório
\end{paraestudar}

\begin{pratique}
  \item Monte manualmente, byte a byte, um quadro de comando e confira com uma captura real.
  \item Adicione um campo a uma mensagem sem quebrar app ou firmware antigos e explique a estratégia.
\end{pratique}

\topico{ESP32 e ESP-IDF}
Estrutura de projeto ESP-IDF (componentes, \code{main}), \code{idf.py}, \code{sdkconfig} e menuconfig, tabela de partições e logs.

\begin{noproduto}
O \code{ergo\_fes\_connect\_ble\_firmware} é um projeto ESP-IDF.
\end{noproduto}

\begin{paraestudar}
  \item \leitura{https://docs.espressif.com/projects/esp-idf/en/latest/esp32/}{Documentação do ESP-IDF, seção Get Started}
\end{paraestudar}

\begin{pratique}
  \item Compile e grave o firmware BLE e acompanhe o boot pelo monitor serial.
  \item Mude uma opção pelo menuconfig e mostre o diff no \code{sdkconfig}.
\end{pratique}

\topico{FreeRTOS}
Tasks e prioridades, escalonador preemptivo, filas, semáforos, mutexes e callbacks. Um RTOS troca a previsibilidade do superloop por concorrência real.

\begin{noproduto}
O firmware BLE roda sobre FreeRTOS, integrado ao ESP-IDF.
\end{noproduto}

\begin{paraestudar}
  \item \leitura{https://www.freertos.org/}{FreeRTOS.org e o livro gratuito Mastering the FreeRTOS Real Time Kernel}
\end{paraestudar}

\begin{pratique}
  \item Liste as tasks do firmware BLE com prioridade e tamanho de stack, e meça o stack livre de cada uma.
  \item Explique como um dado sai da task da UART e chega à pilha BLE sem condição de corrida.
\end{pratique}

\topico{Bluetooth Low Energy: lado periférico}
GAP e advertising, GATT (services, characteristics, descriptor CCC), notify e indicate, MTU. É como um dispositivo se anuncia e expõe dados para um celular.

\begin{noproduto}
O ESP32 do Connect é o periférico BLE e usa a pilha Bluedroid do ESP-IDF.
\end{noproduto}

\begin{paraestudar}
  \item \leitura{https://academy.nordicsemi.com/}{Nordic DevAcademy: Bluetooth Low Energy Fundamentals}
  \item \leitura{https://www.novelbits.io/}{Novel Bits, artigos introdutórios de BLE}
  \item Exemplos gatt\_server do ESP-IDF
\end{paraestudar}

\begin{pratique}
  \item Conecte-se com o nRF Connect, liste services e characteristics do Connect e associe cada UUID à sua função.
  \item Explique a diferença prática entre notify e indicate e onde cada um é usado no projeto.
\end{pratique}

\topico{Ponte entre BLE e UART}
Roteamento de comandos entre characteristics BLE e quadros seriais, validação de toda entrada vinda do celular e limites da tabela de atributos.

\begin{noproduto}
O firmware BLE faz exatamente essa ponte entre o app e o ATSAM.
\end{noproduto}

\begin{paraestudar}
  \item Código da ponte no firmware BLE
\end{paraestudar}

\begin{pratique}
  \item Desenhe o caminho de uma escrita BLE até o quadro UART correspondente.
  \item Envie pelo nRF Connect um payload malformado e verifique se o firmware o rejeita.
\end{pratique}

\topico{Código compartilhado entre plataformas}
Bibliotecas comuns como submódulo e uma camada de abstração de plataforma, para que o mesmo código rode em microcontroladores diferentes e até no PC.

\begin{noproduto}
A pasta \code{Modules/} é compartilhada entre os firmwares, com abstração para ATSAM, ESP e STM32.
\end{noproduto}

\begin{paraestudar}
  \item Pasta Modules/ do repositório
\end{paraestudar}

\begin{pratique}
  \item Identifique o que, em \code{Modules/}, depende de plataforma e como essa dependência é injetada.
  \item Compile um módulo comum para o PC com um stub de plataforma.
\end{pratique}

\topico{Firmware de segurança crítica}
Defesa em profundidade, watchdog, falha segura, interação com as proteções de hardware e diferenças entre builds de depuração e de produção. O firmware nunca é a única barreira, mas precisa ser uma boa barreira.

\begin{noproduto}
No firmware stimulus, a alta tensão só é habilitada em build Release (com \code{NDEBUG}).
\end{noproduto}

\begin{paraestudar}
  \item IEC 62304 (ciclo de vida de software médico)
  \item Livro: Embedded Software Development for Safety-Critical Systems (Chris Hobbs)
\end{paraestudar}

\begin{pratique}
  \item Liste o que impede o firmware de ligar a alta tensão num build Debug e confirme no código.
  \item Simule um travamento do superloop e verifique se o watchdog leva o sistema a um estado seguro.
\end{pratique}

\topico{Testes de firmware}
Testes unitários no PC, testes de bancada, regressão e rastreabilidade entre requisito, código e teste. Firmware testado no PC é firmware que pode mudar sem medo.

\begin{noproduto}
A lógica pura (protocolo, máquinas de estados) pode ser testada no PC; o resto, em bancada.
\end{noproduto}

\begin{paraestudar}
  \item \leitura{https://google.github.io/googletest/}{GoogleTest}
  \item \leitura{https://github.com/ThrowTheSwitch/Unity}{Unity}
  \item \leitura{https://cpputest.github.io/}{CppUTest}
  \item Livro: Test-Driven Development for Embedded C (James Grenning)
\end{paraestudar}

\begin{pratique}
  \item Escreva testes no PC para a serialização de uma mensagem do \code{WeProtocol}.
  \item Escreva um roteiro de teste de bancada para a sequência de início e parada de estímulo.
\end{pratique}

% Parte 3: Software (gerado a partir da apostila NeuroDev; edite à vontade)
\begin{parte}{Software}{Desenvolvimento de aplicativos móveis com Flutter, do Dart básico à comunicação com dispositivos por Bluetooth, backend e lojas de aplicativo. O app do Connect é o caso de estudo do começo ao fim.}
  \item Construir telas, estado e navegação em Flutter com arquitetura em camadas.
  \item Falar com um dispositivo por BLE respeitando o contrato com o firmware.
  \item Testar, integrar com backend e publicar um app de forma confiável.
\end{parte}

\nivel{Básico}{Dart, Flutter, ambiente e projetos Android e iOS.}

\topico{Linguagem Dart}
Tipos, null safety, classes e mixins, \code{async}/\code{await}, \code{Future} e \code{Stream}. Dart é a linguagem do Flutter e tem um modelo assíncrono que aparece em todo lugar.

\begin{noproduto}
Dados que chegam do dispositivo de forma contínua, como bateria e estado, são modelados naturalmente como \code{Stream}.
\end{noproduto}

\begin{paraestudar}
  \item \leitura{https://dart.dev/language}{Dart language tour}
  \item \leitura{https://dart.dev/null-safety}{Sound null safety}
  \item \leitura{https://dartpad.dev/}{DartPad, para testar no navegador}
\end{paraestudar}

\begin{pratique}
  \item Escreva um \code{Stream} que emita a contagem regressiva de um exercício de 10 s e consuma com \code{await for}.
  \item Reescreva uma classe com campos anuláveis para que o compilador garanta a inicialização.
\end{pratique}

\topico{Fundamentos de Flutter}
Widgets, árvore de widgets, Stateless e Stateful, ciclo de vida, layout (Row, Column, Flex) e navegação. Em Flutter, toda a interface é uma árvore de widgets reconstruída a partir do estado.

\begin{noproduto}
Todas as telas do \code{app\_main} são composições de widgets sobre esses fundamentos.
\end{noproduto}

\begin{paraestudar}
  \item \leitura{https://docs.flutter.dev/}{Documentação do Flutter}
  \item Flutterando (comunidade brasileira) no YouTube
\end{paraestudar}

\begin{pratique}
  \item Monte uma tela de controle com slider de amplitude e botão iniciar/parar, usando só estado local.
  \item Explique quando o \code{build} é chamado e como evitar reconstruções desnecessárias.
\end{pratique}

\topico{Ambiente de desenvolvimento}
Flutter SDK, FVM para fixar a versão, \code{pubspec.yaml}, pacotes e Android Studio/Xcode. Um ambiente reprodutível evita o clássico “na minha máquina funciona”.

\begin{noproduto}
O projeto fixa o Flutter 3.19.6 no \code{.fvmrc}.
\end{noproduto}

\begin{paraestudar}
  \item \leitura{https://fvm.app/}{FVM}
  \item \leitura{https://pub.dev/}{pub.dev}
\end{paraestudar}

\begin{pratique}
  \item Configure o FVM, rode o \code{app\_main} num emulador e anote cada problema encontrado (vira FAQ desta apostila).
  \item Explique a diferença entre \code{pubspec.yaml} e \code{pubspec.lock} e qual deles vai para o git.
\end{pratique}

\topico{Projetos Android e iOS}
Estrutura das pastas \code{android/} e \code{ios/}, permissões, assinatura e build de release. Por baixo do Flutter continuam existindo dois projetos nativos.

\begin{noproduto}
Escanear e conectar a dispositivos BLE exige permissões específicas em cada plataforma; confira as que o projeto declara.
\end{noproduto}

\begin{paraestudar}
  \item \leitura{https://docs.flutter.dev/deployment}{Flutter: deployment}
\end{paraestudar}

\begin{pratique}
  \item Liste as permissões de Bluetooth declaradas em Android e iOS e explique por que cada uma existe.
  \item Gere um APK de release assinado com uma chave de teste.
\end{pratique}

\nivel{Intermediário}{Estado, módulos, camadas, pacotes, dados e qualidade.}

\topico{Gerência de estado com BLoC}
\code{flutter\_bloc}: eventos, estados, Bloc e Cubit. Separa a lógica da interface e torna o comportamento previsível e testável.

\begin{noproduto}
O app usa \code{flutter\_bloc}, com Blocs e Cubits.
\end{noproduto}

\begin{paraestudar}
  \item \leitura{https://bloclibrary.dev/}{bloclibrary.dev}
\end{paraestudar}

\begin{pratique}
  \item Refaça a tela do exercício 3.2 com um Cubit.
  \item Desenhe o fluxo evento → estado de um Bloc real do app.
\end{pratique}

\topico{Modularização e injeção de dependência}
Dividir o app em módulos com rotas e dependências próprias, e injetar serviços em vez de criá-los à mão. Facilita trabalho em paralelo e testes.

\begin{noproduto}
O app usa \code{flutter\_modular}, com módulos como \code{fes}, \code{breath}, \code{devices}, \code{patient} e \code{medic}.
\end{noproduto}

\begin{paraestudar}
  \item \leitura{https://modular.flutterando.com.br/}{Documentação do flutter\_modular}
\end{paraestudar}

\begin{pratique}
  \item Mapeie os módulos do app e as rotas de cada um.
  \item Crie um módulo vazio com uma rota e um bind e navegue até ele.
\end{pratique}

\topico{Arquitetura em camadas}
Services, repositories, adapters e helpers organizados em camadas com regras claras de quem pode importar quem. A interface não deve saber como os dados chegam.

\begin{noproduto}
O \code{app\_main} separa \code{core}, \code{domain}, \code{modules} e \code{ui}.
\end{noproduto}

\begin{paraestudar}
  \item Código do app\_main
  \item Livro: Clean Architecture (Robert C. Martin), capítulos sobre dependências
\end{paraestudar}

\begin{pratique}
  \item Siga uma funcionalidade do botão na tela até o service que fala com o BLE e nomeie cada camada.
  \item Encontre uma dependência que viole as camadas, ou mostre que não há.
\end{pratique}

\topico{Pacotes locais e monorepo}
Dependências por \code{path}, submódulos git e reuso de pacotes entre projetos. Permite compartilhar código sem publicar nada.

\begin{noproduto}
O app usa pacotes locais como \code{we\_bluetooth} e \code{we\_exception}.
\end{noproduto}

\begin{paraestudar}
  \item \leitura{https://dart.dev/tools/pub/dependencies}{Dart: package dependencies}
\end{paraestudar}

\begin{pratique}
  \item Altere uma função num pacote local e verifique a mudança no app sem publicar nada.
  \item Explique o que acontece se o submódulo de um pacote estiver num commit diferente do esperado.
\end{pratique}

\topico{Serialização e geração de código}
\code{json\_serializable}, \code{build\_runner} e \code{equatable}. Gerar o código repetitivo evita erros de digitação em centenas de campos.

\begin{noproduto}
O app usa \code{json\_serializable}, \code{build\_runner} e \code{equatable}.
\end{noproduto}

\begin{paraestudar}
  \item \leitura{https://pub.dev/packages/json_serializable}{json\_serializable no pub.dev}
  \item \leitura{https://pub.dev/packages/equatable}{equatable no pub.dev}
\end{paraestudar}

\begin{pratique}
  \item Crie um modelo com \code{json\_serializable}, gere o código e escreva um teste de ida e volta.
  \item Explique por que estados do BLoC costumam estender \code{Equatable}.
\end{pratique}

\topico{Internacionalização}
Textos traduzidos, datas, números e formatos por idioma. Pensar em idiomas desde o início custa pouco; adaptar depois custa caro.

\begin{noproduto}
O app usa \code{slang} e \code{intl}.
\end{noproduto}

\begin{paraestudar}
  \item \leitura{https://pub.dev/packages/slang}{slang no pub.dev}
  \item \leitura{https://pub.dev/packages/intl}{intl no pub.dev}
\end{paraestudar}

\begin{pratique}
  \item Adicione uma string nova em português e inglês e use-a numa tela.
  \item Formate a data de uma sessão nos dois idiomas.
\end{pratique}

\topico{Persistência local}
Banco local, armazenamento seguro e caminhos de arquivos. Dados de sessão e credenciais exigem cuidados diferentes.

\begin{noproduto}
O app usa Hive, armazenamento seguro de sessão e \code{path\_provider}.
\end{noproduto}

\begin{paraestudar}
  \item \leitura{https://pub.dev/packages/hive}{hive no pub.dev}
  \item \leitura{https://pub.dev/packages/path_provider}{path\_provider no pub.dev}
\end{paraestudar}

\begin{pratique}
  \item Liste o que o app guarda localmente e onde.
  \item Explique o que não deve ir para armazenamento comum e por quê.
\end{pratique}

\topico{Formulários e validação}
Formulários, máscaras de entrada e validação de dados como CPF. Validar na entrada é mais barato que corrigir no banco.

\begin{noproduto}
O app tem máscaras, validação de CPF e helpers de formato.
\end{noproduto}

\begin{paraestudar}
  \item \leitura{https://docs.flutter.dev/cookbook/forms/validation}{Flutter: form validation}
\end{paraestudar}

\begin{pratique}
  \item Implemente a validação de CPF pelos dígitos verificadores e escreva testes com casos válidos e inválidos.
  \item Adicione uma máscara de telefone a um campo existente.
\end{pratique}

\topico{Lint e qualidade}
\code{analysis\_options.yaml}, regras de lint e \code{dart format}. O analisador estático pega erros antes de o app rodar.

\begin{noproduto}
O projeto usa as regras do \code{flutter\_lints}.
\end{noproduto}

\begin{paraestudar}
  \item \leitura{https://dart.dev/tools/analysis}{Dart: customizing static analysis}
\end{paraestudar}

\begin{pratique}
  \item Rode \code{flutter analyze} e corrija três avisos.
  \item Proponha uma regra de lint a mais para o projeto e justifique.
\end{pratique}

\nivel{Avançado}{BLE, bytes, backend, fluxos clínicos, testes e publicação.}

\topico{Bluetooth Low Energy: lado central}
Scan, conexão, descoberta de serviços, leitura, escrita e notify, permissões e diferenças entre Android e iOS. É o lado do celular na conversa BLE.

\begin{noproduto}
O app usa \code{flutter\_reactive\_ble}, junto com o pacote local \code{bluetooth}.
\end{noproduto}

\begin{paraestudar}
  \item \leitura{https://pub.dev/packages/flutter_reactive_ble}{flutter\_reactive\_ble no pub.dev}
  \item \leitura{https://academy.nordicsemi.com/}{Nordic DevAcademy: Bluetooth Low Energy Fundamentals}
\end{paraestudar}

\begin{pratique}
  \item Escreva um app mínimo que encontre o Connect, conecte e mostre o nível de bateria em tempo real.
  \item Trate a desconexão no meio de uma sessão e defina o que a interface deve mostrar.
\end{pratique}

\topico{Conversão de bytes}
Serialização binária com \code{ByteData}, endianness e o espelho, no app, das structs definidas no firmware. Um byte fora do lugar muda o valor inteiro.

\begin{noproduto}
O pacote \code{byte\_conversion} faz essa ponte com o firmware.
\end{noproduto}

\begin{paraestudar}
  \item \leitura{https://api.dart.dev/stable/dart-typed_data/ByteData-class.html}{ByteData (dart:typed\_data)}
  \item Pacote byte\_conversion do repositório
\end{paraestudar}

\begin{pratique}
  \item Converta uma struct do firmware (com \code{int16}, \code{uint32} e \code{float}) para bytes e de volta, com teste.
  \item Mostre o bug que aparece quando app e firmware discordam do endianness.
\end{pratique}

\topico{Contrato entre app e firmware}
UUIDs, layout de cada característica, versionamento de protocolo e como evoluir sem quebrar dispositivos ou apps antigos. O contrato é o documento mais importante entre duas equipes.

\begin{noproduto}
O contrato do Connect está no documento técnico e no pacote \code{bluetooth}.
\end{noproduto}

\begin{paraestudar}
  \item Documento técnico do Connect e código do pacote bluetooth
\end{paraestudar}

\begin{pratique}
  \item Monte uma tabela característica × layout de bytes × quem escreve × quem lê.
  \item Proponha como o app deve se comportar ao encontrar um firmware com versão de protocolo mais nova.
\end{pratique}

\topico{Concorrência no app}
Streams, locks assíncronos e controle de comandos simultâneos, para que uma escrita não atropele outra.

\begin{noproduto}
O app usa o pacote \code{synchronized} para serializar comandos BLE.
\end{noproduto}

\begin{paraestudar}
  \item \leitura{https://pub.dev/packages/synchronized}{synchronized no pub.dev}
  \item \leitura{https://dart.dev/libraries/async/async-await}{Dart: asynchronous programming}
\end{paraestudar}

\begin{pratique}
  \item Dispare duas escritas BLE ao mesmo tempo sem lock e observe o resultado; depois proteja com \code{synchronized}.
  \item Explique a diferença entre isolates e código assíncrono no mesmo isolate.
\end{pratique}

\topico{Comunicação com backend}
HTTP (interceptors, timeouts), gRPC, autenticação JWT e upload de arquivos. É como o app guarda e busca dados fora do celular.

\begin{noproduto}
O app usa Dio para HTTP, gRPC e autenticação JWT.
\end{noproduto}

\begin{paraestudar}
  \item \leitura{https://pub.dev/packages/dio}{dio no pub.dev}
  \item \leitura{https://grpc.io/docs/languages/dart/}{gRPC para Dart}
  \item \leitura{https://jwt.io/introduction}{Introdução a JWT}
\end{paraestudar}

\begin{pratique}
  \item Mapeie as chamadas do app ao backend e quais usam HTTP e quais usam gRPC.
  \item Explique como o app renova um token expirado sem deslogar o usuário.
\end{pratique}

\topico{Firebase}
Autenticação, banco de documentos, notificações push e relatórios de falha como serviço.

\begin{noproduto}
O app usa Firebase Auth, Firestore, Cloud Messaging e Crashlytics.
\end{noproduto}

\begin{paraestudar}
  \item \leitura{https://firebase.google.com/docs/flutter/setup}{Firebase para Flutter}
\end{paraestudar}

\begin{pratique}
  \item Liste o que o app usa de cada serviço Firebase.
  \item Force um erro não tratado num build de teste e encontre-o no Crashlytics.
\end{pratique}

\topico{Tratamento de erros padronizado}
Códigos de erro únicos e mapeamento de falhas de rede, de protocolo e de dispositivo para mensagens que o usuário entende.

\begin{noproduto}
O pacote \code{we\_exception} padroniza os erros e mapeia respostas HTTP e gRPC.
\end{noproduto}

\begin{paraestudar}
  \item Código do pacote we\_exception
\end{paraestudar}

\begin{pratique}
  \item Siga um erro 401 do backend até a mensagem exibida na tela.
  \item Adicione um código de erro para falha de conexão BLE e use-o em um ponto do app.
\end{pratique}

\topico{Fluxos clínicos no app}
Pré e pós-sessão, execução e interrupção de exercício, relatório e agenda. Em um app de saúde, o fluxo do usuário é parte da segurança.

\begin{noproduto}
Os módulos \code{fes}, \code{breath}, \code{patient} e \code{medic} implementam esses fluxos, incluindo o chat de pré e pós- sessão.
\end{noproduto}

\begin{paraestudar}
  \item Código dos módulos fes, breath, patient e medic
\end{paraestudar}

\begin{pratique}
  \item Desenhe a sequência de telas de uma sessão completa, do login ao relatório.
  \item Liste todos os caminhos pelos quais uma sessão pode ser interrompida e o que acontece no dispositivo em cada um.
\end{pratique}

\topico{Gráficos e visualização}
Gráficos de sessão e de evolução, e indicadores de progresso. Bons gráficos ajudam o profissional a decidir.

\begin{noproduto}
O app usa Syncfusion Charts.
\end{noproduto}

\begin{paraestudar}
  \item \leitura{https://pub.dev/packages/syncfusion_flutter_charts}{syncfusion\_flutter\_charts no pub.dev}
\end{paraestudar}

\begin{pratique}
  \item Crie um gráfico de linha com a amplitude ao longo de uma sessão simulada.
  \item Verifique a licença do Syncfusion usada pelo projeto e suas condições.
\end{pratique}

\topico{Testes}
Testes unitários, de widget e de integração, e mocks da camada de hardware para testar sem dispositivo.

\begin{noproduto}
Com um mock do serviço BLE, a maior parte do app pode ser testada sem o Connect por perto.
\end{noproduto}

\begin{paraestudar}
  \item \leitura{https://docs.flutter.dev/testing/overview}{Flutter: testing}
  \item \leitura{https://pub.dev/packages/mocktail}{mocktail no pub.dev}
\end{paraestudar}

\begin{pratique}
  \item Escreva um teste de Cubit para o fluxo iniciar → executando → parar.
  \item Crie um mock do serviço BLE e teste a tela de controle sem dispositivo.
\end{pratique}

\topico{Distribuição}
Builds de release, Play Store e App Store, trilhas de teste e versionamento.

\begin{noproduto}
A versão do app é definida no campo \code{version} do \code{pubspec.yaml}.
\end{noproduto}

\begin{paraestudar}
  \item \leitura{https://docs.flutter.dev/deployment}{Flutter: deployment}
\end{paraestudar}

\begin{pratique}
  \item Descreva o processo de release do app passo a passo (vira checklist da equipe).
  \item Explique a diferença entre \code{versionName} e \code{versionCode} e como o pubspec define os dois.
\end{pratique}

% Parte 4: Integração do sistema
\begin{parte}{Integração do sistema}{A trilha final, comum a todas as áreas. Aqui o membro deixa de olhar só para a sua camada e passa a responder pelo sistema inteiro: do toque na tela até a corrente no eletrodo.}
  \item Seguir um comando por todas as camadas e localizar onde um problema nasce.
  \item Implementar uma funcionalidade ponta a ponta, do firmware ao app.
  \item Documentar, analisar riscos e validar um sistema completo.
\end{parte}

\topico{O caminho completo de um comando}
Todo sistema embarcado conectado tem uma cadeia de ida (comando) e de volta (telemetria). Conhecer cada elo, sua latência e seus pontos de falha é o que separa quem conserta sintomas de quem resolve causas.

\begin{noproduto}
Ida: app → BLE → ESP32 → UART → ATSAM → DAC → fonte de corrente → eletrodos. Volta: bateria, impedância e estado fazem o caminho inverso.
\end{noproduto}

\begin{paraestudar}
  \item Documento técnico do Connect
  \item Tópicos 2.21, 2.22, 2.26, 3.14 e 3.16 deste documento
\end{paraestudar}

\begin{pratique}
  \item Desenhe um diagrama de sequência (Mermaid) do comando “iniciar estímulo”, com todas as mensagens e camadas.
  \item Estime a latência de cada trecho e meça a latência total do toque ao primeiro pulso.
\end{pratique}

\topico{Uma funcionalidade ponta a ponta}
Adicionar uma funcionalidade a um produto com várias camadas exige mudanças coordenadas e compatíveis em cada uma. Planejar antes de codar evita retrabalho.

\begin{noproduto}
Estudo de caso: o GVS, que exigiu quatro passos no ATSAM, uma característica nova no ESP32 e suporte no app.
\end{noproduto}

\begin{paraestudar}
  \item Histórico de commits e PRs da implementação do GVS nos três repositórios
\end{paraestudar}

\begin{pratique}
  \item Reconstrua, a partir dos commits, a lista de mudanças que o GVS exigiu em cada repositório.
  \item Escreva o plano de implementação de uma funcionalidade simples (por exemplo, um contador de sessões) nas três camadas.
\end{pratique}

\topico{Depuração entre camadas}
Sniffer BLE, analisador lógico e logs com horário correlacionados. A técnica é sempre a mesma: observar a fronteira entre duas camadas para descobrir de qual lado o problema está.

\begin{noproduto}
As fronteiras observáveis do Connect são o enlace BLE (nRF Connect) e a UART entre ESP32 e ATSAM (analisador lógico).
\end{noproduto}

\begin{paraestudar}
  \item \leitura{https://sigrok.org/}{sigrok / PulseView}
  \item \leitura{https://www.saleae.com/}{Saleae Logic}
  \item App nRF Connect for Mobile (Nordic)
\end{paraestudar}

\begin{pratique}
  \item Capture ao mesmo tempo o tráfego BLE e a UART durante um comando e alinhe os eventos numa linha do tempo.
  \item Introduza um atraso artificial no firmware BLE e localize-o só pelas capturas.
\end{pratique}

\topico{Documentação de engenharia}
Documento técnico e diagramas de arquitetura, sequência e estados, mantidos junto do código. Documentação boa é a que o próximo membro consegue usar sem perguntar.

\begin{noproduto}
Módulos novos, como o de tDCS, devem terminar com um documento técnico próprio.
\end{noproduto}

\begin{paraestudar}
  \item \leitura{https://mermaid.js.org/}{Mermaid}
  \item \leitura{https://c4model.com/}{C4 model, para diagramas de arquitetura}
\end{paraestudar}

\begin{pratique}
  \item Escreva a documentação técnica de um módulo que você estudou, com ao menos um diagrama.
  \item Revise a documentação de um colega e aponte o que não daria para entender sem ler o código.
\end{pratique}

\topico{Análise de segurança do sistema}
FMEA, árvore de falhas e o mapeamento de cada mitigação entre hardware, firmware e app. Em equipamento médico, todo risco relevante precisa de pelo menos uma barreira registrada.

\begin{noproduto}
Um bom exercício é perda de BLE durante o estímulo: o que o app, o ESP32, o ATSAM e o hardware fazem.
\end{noproduto}

\begin{paraestudar}
  \item ISO 14971 (gestão de risco)
  \item Material de FMEA de disciplinas de confiabilidade ou da equipe
\end{paraestudar}

\begin{pratique}
  \item Faça o FMEA de uma falha com causa, efeito, severidade e mitigação em cada camada.
  \item Monte a árvore de falhas do evento “corrente acima do limite no paciente”.
\end{pratique}

\topico{Verificação e validação}
Testes de bancada, testes ponta a ponta, regressão e rastreabilidade de requisito até teste. Verificar é provar que se construiu certo; validar é provar que se construiu a coisa certa.

\begin{noproduto}
Um roteiro de regressão ponta a ponta antes de cada release protege as três bases de código ao mesmo tempo.
\end{noproduto}

\begin{paraestudar}
  \item IEC 62304 (verificação de software médico)
  \item Tópicos 1.19, 2.29 e 3.23 deste documento
\end{paraestudar}

\begin{pratique}
  \item Monte uma matriz de rastreabilidade para três requisitos: requisito → implementação → teste → resultado.
  \item Escreva um roteiro de regressão ponta a ponta que possa ser rodado antes de cada release.
\end{pratique}


%% ------------------------------------------------------------
%% APÊNDICES
%% ------------------------------------------------------------
\apendice{Glossário}
Termos que aparecem ao longo da apostila, em ordem alfabética.
\par\vspace{2mm}
\begin{glossario}
  \termo{ADC / DAC}{Conversor analógico-digital e digital-analógico: traduzem tensões em números e números em tensões.}
\termo{BLE}{Bluetooth Low Energy, versão de baixo consumo do Bluetooth usada entre o celular e o dispositivo.}
\termo{BLoC / Cubit}{Padrão de gerência de estado no Flutter que separa a lógica da interface.}
\termo{Bootstrap}{Técnica de gate driver que usa um capacitor carregado para acionar o transistor do lado alto de uma ponte.}
\termo{CCC}{Client Characteristic Configuration: descritor BLE que liga ou desliga notify e indicate.}
\termo{Checksum}{Valor calculado sobre um quadro de dados para detectar erros de transmissão.}
\termo{Compliance}{Maior tensão que uma fonte de corrente consegue aplicar mantendo a corrente programada.}
\termo{Debounce}{Filtragem do ruído mecânico de botões e chaves.}
\termo{Endianness}{Ordem dos bytes de um número na memória ou num protocolo (little-endian ou big-endian).}
\termo{ESR}{Resistência série equivalente de um capacitor; define perdas e ripple.}
\termo{FES}{Estimulação elétrica funcional: pulsos que provocam contração muscular útil.}
\termo{FMEA}{Análise de modos de falha e efeitos: tabela de falhas possíveis, consequências e mitigações.}
\termo{GATT}{Estrutura de services e characteristics pela qual um periférico BLE expõe dados.}
\termo{GVS}{Estimulação vestibular galvânica: corrente aplicada atrás das orelhas que atua no equilíbrio.}
\termo{ISR}{Rotina de serviço de interrupção, código executado quando um evento de hardware acontece.}
\termo{Jitter}{Variação indesejada no tempo de um evento que deveria ser periódico.}
\termo{Latch}{Circuito que mantém um estado (por exemplo, falha) até ser rearmado.}
\termo{MTU}{Maior pacote de dados aceito num enlace; no BLE define quantos bytes cabem numa operação.}
\termo{NVIC}{Controlador de interrupções dos núcleos ARM Cortex-M.}
\termo{Open-drain}{Saída que só puxa para o terra; precisa de pull-up para gerar nível alto.}
\termo{RAII}{Técnica de C++ que amarra a aquisição e a liberação de um recurso ao ciclo de vida de um objeto.}
\termo{SERCOM}{Periférico serial configurável do ATSAMC20 (SPI, I²C ou USART).}
\termo{Shunt}{Resistor de baixo valor usado para medir corrente pela queda de tensão.}
\termo{SWD}{Serial Wire Debug: interface de gravação e depuração dos ARM Cortex-M.}
\termo{tDCS}{Estimulação transcraniana por corrente contínua, de baixa intensidade, aplicada no couro cabeludo.}
\termo{UART}{Interface serial assíncrona de dois fios (TX e RX).}
\termo{Watchdog}{Temporizador que reinicia o sistema se o firmware parar de sinalizar que está vivo.}

\end{glossario}

\apendice{Ficha de progresso}
Marque cada tópico concluído e anote a data. A ficha também ajuda quem acompanha sua formação a saber onde você está.
\par\vspace{2mm}
\fichaprogresso

\apendice{Como contribuir}
A NeuroDev é um documento vivo. Sempre que estudar um tópico, deixe-o um pouco melhor: troque um link quebrado, acrescente um vídeo que ajudou, proponha um exercício ou corrija algo que mudou no código.

\begin{enumerate}
  \item Faça a mudança no documento-fonte da apostila, nunca só no PDF.
  \item Siga o modelo de tópico abaixo, mantendo a numeração da trilha.
  \item Confira se o quadro No Connect ainda corresponde ao código atual.
  \item Registre a mudança no histórico de revisões e peça a revisão de um colega.
\end{enumerate}

\secao{Modelo de tópico}
\begingroup\naocontarexercicios
\begin{tcolorbox}[enhanced, colback=white, colframe=nro-pia, boxrule=0.5pt, arc=0pt, left=5mm, right=5mm, top=4mm, bottom=3mm]
  {\fonttitulo\fontsize{14}{17}\selectfont\color{nro-axon}X.Y\hspace{2mm}\color{nro-tinta}Nome do tópico}\par\vspace{1mm}
  Texto de abertura: o conceito geral, em 2 a 3 frases.
  \begin{noproduto}
    Onde o conceito aparece nos repositórios ou esquemáticos.
  \end{noproduto}
  \begin{paraestudar}
    \item Link, livro ou vídeo (com capítulo, se fizer sentido)
    \item Documento ou código interno relacionado
  \end{paraestudar}
  \begin{pratique}
    \item Exercício conceitual
    \item Exercício prático ligado ao código ou ao hardware real
  \end{pratique}
\end{tcolorbox}
\endgroup

\secao{Controle do documento}
\controledocumento

\secao{Histórico de revisões}
\begin{historico}
  \linharevisao{25/10/2026}{A}{Gabriel Matos Veneziano}{Primeira edição: 5 partes, 85 tópicos, glossário e ficha de progresso.}
\end{historico}

\contracapa{Apostila de formação técnica da equipe NeuroDynamics. Encontrou um erro ou tem uma sugestão? Veja \enquote{Como contribuir} e ajude com o desenvolvimento da apostila.}

\end{document}
